\section*{Plagiarism and AI declaration}
\phantomsection  % gives \label a target: \section* sets no anchor of its own
\label{sec:plagiarism}
% TODO
During the course of this project, I made use of Claude Code to assist me in the research process. It was used in general cases where redundant tasks / code was needed (soft use) and then cases where a citation for the tool is required (hard use) in accordance to the guidelines set out in "Guidelines for using GenAI tools in Learning and Assessment". Listed below are all the specific items that I used the tool for:
% Numbered so the report can point at an item: \ref{ai:eq-cue} prints "AI6a" and links here.
% Write (\ref{ai:...}) in the text, the way \eqref writes (3).
\begin{enumerate}
    \renewcommand{\theenumi}{AI\arabic{enumi}}
    \renewcommand{\theenumii}{\alph{enumii}}
    \item \label{ai:research} \textbf{Research:} I prompted Claude Code to assist me in finding relevant research papers and articles that may have potentially solved the TSE task. I used the tool to highlight short comings and advantages of different approaches. It was from this research that I decided on the BSRNN model. Specifically for the extension of the BSRNN model, I prompted the tool to analyse the short comings of the model and to suggest potential avenues to explore to improve the model. I then decided on the specific design choices for the extension from my own understanding of the tool's response.
    \item \label{ai:writing} \textbf{Writing assistant:} I wrote the report myself and made use of Claude Code in fixing grammar, spelling and sentence structure. The tool was not given access to writing directly to the report, but rather would suggest corrections if a claim or sentence was not clear or inaccurate. I would ask the tool if my summary of a paper or method was fair and accurate to not misrepresent the work of others or my own work. The formulations and equations for the described methods were sourced by the original authors, but I would first rewrite them from my understanding and then ask the tool to help change the notation to be consistent with the rest of my report.
    \item \label{ai:diagrams} \textbf{Diagrams:} Figure~\ref{fig:baseline-architecture} specifically I prompted Claude Code to assist with the TikZ code for the best layout of the diagram. The other figures in the report were reformatted and editted by Claude Code to have a scientific style and to be consistent with the rest of the report.
    \item \label{ai:loss} \textbf{Loss function:} The terms in the objective function (Section~\ref{sec:objective}) were sourced from CARTSE \cite{li2026cartse} and Yu et al. \cite{yu2023high} where explicitly stated. 
    \item \label{ai:code-soft} \textbf{Code (soft):} I made use of Claude Code to assist me in documenting the decisions from my discussion with the tool into the code and recorded them seperately. All decisions were logged into the documentation and can be found in the ```docs/decisions''' directory in the repository found at \url{https://github.com/grantxxcoder/TargetSpeakerExtractor}. For the code pertinent to the baseline and extension model for the actual TSE task, I made use of Claude Code to debug and fix errors. Furthermore, for the baseline, because it needed to follow the architecture found in the REAL TSE paper, I needed assistance in implementing the architecture as described in the paper. I prompted the tool to verify that the implementation was consistent with the paper and to suggest changes if it was not. For function specific code, I marked all the functions that I used Claude Code to debug and fix errors with a reference.
    \item \label{ai:code-hard} \textbf{Code (hard):} I orginally wrote my code in scripts but Kaggle only accepts the files in a notebook format. Thus I prompted the tool to fully convert my scripts into a bundle of configuration files and notebooks that can be run on Kaggle. The tool then created multiple hyperparameter configuration files (.\textit{yaml}), notebooks for training different models and dataset loaders. I then prompted the tool to verify that the notebooks were consistent with the scripts and that they would run on Kaggle. The generated files only contain the code required to train and save the models and no additional logic is used to change the model or training procedure. For the files and functions coded that were completely generated by the tool, I expressely marked in the repository as Claude Code (Hard) generated. These are files and functions that do not form part of the core logic, but rather for setup, configuration, tests, or redundant tasks.
    \item \label{ai:equations} \textbf{Equations and Formulations:}
    \begin{enumerate}
        \item \label{ai:eq-cue} For the first extension, the splitting the cue into parts idea came after the initial diagnosis of the failures of the BSRNN. I prompted Claude Code with ideas for improving the speaker representation in which the splitting of the cue into parts was suggested. I used Claude to verify the equations and to ensure consistent notation.
        \item \label{ai:eq-gate} I had the idea of using a two state target present markov model as a conditioning gate for the BSRNN model. I was able to derive the expressions by hand until the KL term myself. For the KL term, Claude Code pointed out that under the mean field assumption the expectation over the neighbouring frames splits into the product of their probabilities and it cited this identity to the paper by Ghahramani and Jordan \cite{ghahramani1997factorial}.  I then worked through each step myself. For the training and inference, Claude Code suggested and derived the equations for the Monte Carlo sampling and the binary Concrete distribution.
    \end{enumerate}
    \item \label{ai:metrics} \textbf{Created Metrics:} The new metrics, the fabrication rate and the interferer rate, were my own ideas, in which I only used Claude to help formulate the equations.
\end{enumerate}
